\documentclass[a4paper,10pt]{article}
\usepackage[en-us,
    % italicdiff,               % switches the differential symbol \(\dd\) to an italic form.
    % draft                     % enables draft mode.
]{YTMstyle}

\title{Example Document}
\author{Tianming You}
\date{\today}

\begin{document}
\maketitle

\tableofcontents

\newpage

\section{Language}
This template currently supports two languages: \texttt{zh-cn} and \texttt{en-us}.

\section{Math Theorems}
\begin{definition}
    This is some definition.
\end{definition}
\begin{theorem}\label{thm:thm1}
    This is a theorem.
\end{theorem}
\begin{remark}
    All \texttt{amsthm} environments share the same counter.
\end{remark}
\begin{example}
    All remarks and examples automatically end with \texttt{\textbackslash diamond}.
\end{example}
\begin{cthdefinition}{Colored Definition}
    This is a colored definition.
\end{cthdefinition}
\begin{cththeorem}{Colored Theorem}
    This is a colored theorem. Colored environments do not share the same counter.
\end{cththeorem}

\section{Options}
\begin{enumerate}
    \item \texttt{en-us} or \texttt{zh-cn}: It specifies the language the article uses. Default: \texttt{en-us}.
    \item \texttt{italicdiff}: It switches the differential symbol \(\dd\) to an italic form \(d\). Default: \texttt{None}.
    \item \texttt{draft}: It enables draft mode, which uses the \texttt{showkeys} package. Default: \texttt{None}.
\end{enumerate}

\section{Packages Included}
\begin{itemize}
    \item Language and encoding: \texttt{ctex} (if \texttt{zh-cn} is used);
    \item Page layout: \texttt{geometry};
    \item Header and footer: \texttt{fancyhdr}, \texttt{lastpage};
    \item Math packages: \texttt{amsmath}, \texttt{amssymb}, \texttt{amsfonts}, \texttt{amsthm}, \texttt{mathrsfs}, \texttt{physics};
    \item Formatting: \texttt{coloredtheorem}, \texttt{enumitem}, \texttt{hyperref}, \texttt{cleveref};
    \item Draft (if \texttt{draft} is used): \texttt{showkeys}; 
    \item Graphics: \texttt{graphicx}, \texttt{quiver}.
\end{itemize}

\section{Other Commands}
\begin{enumerate}
    \item \texttt{\textbackslash ZZ}: \(\mathbb Z\). \texttt{\textbackslash QQ}: \(\mathbb Q\). \texttt{\textbackslash RR}: \(\mathbb R\). \texttt{\textbackslash CC}: \(\mathbb C\).
    \item All commands from packages included.
\end{enumerate}


\end{document}